\frfilename{mt212.tex}
\versiondate{10.9.04}

\def\chaptername{Taksonomi ruang ukur}
\def\sectionname{Ruang lengkap}

\newsection{212}

Dalam dua bagian berikutnya dari bab ini saya memberikan uraian singkat
mengenai teori ruang ukur yang memiliki beberapa sifat yang dijelaskan
dalam \S211. Saya mulai dengan `kelengkapan'. Saya memberikan
fakta-fakta elementer mengenai ruang ukur lengkap dalam 212A-212B;
kemudian saya beralih ke gagasan `pelengkapan' suatu ukuran (212C) dan
hubungannya dengan konsep-konsep teori ukuran lainnya yang telah
diperkenalkan sejauh ini (212D-212G).

\leader{212A}{Proposisi} Setiap ruang ukur yang dikonstruksi
dengan metode \Caratheodory adalah lengkap.

\proof{ Ingat bahwa `metode \Caratheodory' dimulai dari suatu
ukuran luar sembarang $\theta:\Cal PX\to[0,\infty]$ dan menetapkan

\Centerline{$\Sigma=\{E:E\subseteq X,\,\theta A
=\theta(A\cap E)+\theta(A\setminus E)$ untuk setiap $A\subseteq X\}$,
\quad$\mu=\theta\restr\Sigma$}

\noindent (113C). Dalam hal ini, jika $B\subseteq E\in\Sigma$ dan
$\mu E=0$, maka $\theta B=\theta E=0$ (113A(ii)), sehingga untuk setiap
$A\subseteq X$ kita mempunyai

\Centerline{$\theta(A\cap B)+\theta(A\setminus B)
=\theta(A\setminus B)\le\theta A
\le\theta(A\cap B)+\theta(A\setminus B)$,}

\noindent dan $B\in\Sigma$.
}%end of proof of 212A

\leader{212B}{Proposisi} (a) Jika $(X,\Sigma,\mu)$ adalah ruang ukur
lengkap, maka setiap subhimpunan koterabaikan dari $X$ adalah terukur.

(b) Misalkan $(X,\Sigma,\mu)$ suatu ruang ukur lengkap, dan $f$ suatu
fungsi bernilai $[-\infty,\infty]$ yang didefinisikan pada suatu
subhimpunan $X$. Jika $f$ terukur secara virtual\cmmnt{ (yaitu, terdapat
suatu himpunan koterabaikan $E\subseteq X$ sedemikian sehingga
$f\restr E$ terukur)}, maka $f$ terukur.

(c) Misalkan $(X,\Sigma,\mu)$ suatu ruang ukur lengkap, dan $f$ suatu
fungsi bernilai real yang didefinisikan pada suatu subhimpunan
koterabaikan dari $X$. Maka pernyataan-pernyataan berikut ekuivalen,
yakni, jika salah satunya benar maka yang lainnya juga benar:

\quad (i) $f$ terintegralkan;

\quad(ii) $f$ terukur dan $|f|$ terintegralkan;

\quad(iii) $f$ terukur dan terdapat fungsi terintegralkan $g$
sedemikian sehingga $|f|\leae g$.

\proof{{\bf (a)} Jika $E$ koterabaikan, maka $X\setminus E$
terabaikan, sehingga terukur, dan $E$ terukur.

\medskip

{\bf (b)} Misalkan $a\in\Bbb R$. Maka terdapat $H\in\Sigma$ sedemikian
sehingga


\Centerline{$\{x:(f\restr E)(x)\le a\}=H\cap\dom(f\restr E)
=H\cap E\cap\dom f$.}

\noindent Sekarang $F=\{x:x\in\dom f\setminus E,\,f(x)\le a\}$ adalah
subhimpunan dari himpunan terabaikan $X\setminus E$, sehingga terukur,
dan

\Centerline{$\{x:f(x)\le a\}=(F\cup H)\cap\dom f\in\Sigma_{\dom f}$,}

\noindent dengan menuliskan $\Sigma_D=\{D\cap E:E\in\Sigma\}$,
seperti dalam 121A. Karena $a$ sembarang, $f$ terukur (135E).

\medskip

{\bf (c)(i)$\Rightarrow$(ii)} Jika $f$ terintegralkan, maka menurut
122P, $f$ terukur secara virtual dan menurut 122Re, $|f|$
terintegralkan. Menurut (b) di sini, $f$ terukur, sehingga (ii) benar.

\medskip

\quad{\bf (ii)$\Rightarrow$(iii)} adalah langsung.

\medskip

\quad{\bf (iii)$\Rightarrow$(i)} Jika $f$ terukur dan $g$
terintegralkan serta $|f|\leae g$, maka $f$ terukur secara virtual,
$|g|$ terintegralkan, dan $|f|\leae|g|$, sehingga 122P memberi tahu
kita bahwa $f$ terintegralkan.
}%end of proof of 212B

\leader{212C}{Pelengkapan suatu ukuran} Misalkan $(X,\Sigma,\mu)$
sembarang ruang ukur.

\header{212Ca}{\bf (a)} Misalkan $\hat\Sigma$ adalah keluarga semua
himpunan $E\subseteq X$ sedemikian sehingga terdapat $E'$,
$E''\in\Sigma$ dengan $E'\subseteq E\subseteq E''$ dan
$\mu(E''\setminus E')=0$. Maka $\hat\Sigma$ adalah aljabar-sigma dari
subhimpunan-subhimpunan $X$. \prooflet{\Prf\
{(i)} Tentu saja $\emptyset$ termasuk dalam $\hat\Sigma$, karena kita
dapat mengambil $E'=E''=\emptyset$. {(ii)} Jika $E\in\hat\Sigma$,
ambil $E'$, $E''\in\Sigma$ sedemikian sehingga
$E'\subseteq E\subseteq E''$ dan $\mu(E''\setminus E')=0$.
Maka

\Centerline{$X\setminus E''\subseteq X\setminus E
\subseteq X\setminus E'$,
\quad $\mu((X\setminus E')\setminus(X\setminus E''))
=\mu(E''\setminus E')=0$,}

\noindent sehingga $X\setminus E\in\hat\Sigma$. {(iii)} Jika
$\sequencen{E_n}$ adalah suatu barisan dalam $\hat\Sigma$, maka untuk
setiap $n$ pilih $E'_n$, $E_n''\in\Sigma$ sedemikian sehingga
$E'_n\subseteq E_n\subseteq E_n''$ dan
$\mu(E_n''\setminus E'_n)=0$. Tetapkan

\Centerline{$E=\bigcup_{n\in\Bbb N}E_n$,
\quad$E'=\bigcup_{n\in\Bbb N}E'_n$,
\quad$E''=\bigcup_{n\in\Bbb N}E_n''$;}

\noindent maka $E'\subseteq E\subseteq E''$ dan
$E''\setminus E'\subseteq\bigcup_{n\in\Bbb N}(E_n''\setminus E'_n)$
terabaikan, sehingga $E\in\hat\Sigma$. \Qed}

\spheader 212Cb Untuk $E\in\hat\Sigma$, tetapkan

\Centerline{$\hat\mu E=\mu^*E
=\min\{\mu F:E\subseteq F\in\Sigma\}$\dvro{.}{}}

\cmmnt{
\noindent (132A). Patut segera dikemukakan bahwa jika
$E\in\hat\Sigma$, $E'$, $E''\in\Sigma$,
$E'\subseteq E\subseteq E''$, dan $\mu(E''\setminus E')=0$, maka
$\mu E'=\hat\mu E=\mu E''$; hal ini karena

\Centerline{$\mu E'=\mu^*E'\le\mu^*E\le\mu^*E''=\mu E''
=\mu E'+\mu(E''\setminus E')=\mu E'$}

\noindent (dengan mengingat kembali dari 132A, atau mengamati sekarang,
bahwa $\mu^*A\le\mu^*B$ bilamana $A\subseteq B\subseteq X$, dan bahwa
$\mu^*$ sama dengan $\mu$ pada $\Sigma$).
}%end of comment

\header{212Cc}{\bf (c)}\cmmnt{ Sekarang kita mendapati bahwa}
$(X,\hat\Sigma,\hat\mu)$ adalah suatu ruang ukur. \prooflet{\Prf\ (i)
Tentu saja $\hat\mu$, seperti $\mu$, mengambil nilai dalam
$[0,\infty]$. (ii) $\hat\mu\emptyset=\mu\emptyset=0$.
(iii) Misalkan $\sequencen{E_n}$ suatu barisan saling lepas dalam
$\hat\Sigma$, dengan gabungan $E$. Untuk setiap $n\in\Bbb N$ pilih
$E_n'$, $E_n''\in\Sigma$ sedemikian sehingga
$E_n'\subseteq E_n\subseteq E_n''$ dan
$\mu(E_n''\setminus E_n')=0$. Tetapkan
$E'=\bigcup_{n\in\Bbb N}E_n'$,
$E''=\bigcup_{n\in\Bbb N}E_n''$. Maka (seperti dalam (a-iii) di atas)
$E'\subseteq E\subseteq E''$ dan $\mu(E''\setminus E')=0$, sehingga

\Centerline{$\hat\mu E=\mu E'=\sum_{n=0}^{\infty}\mu E_n'
=\sum_{n=0}^{\infty}\hat\mu E_n$}

\noindent karena $\sequencen{E_n'}$, seperti $\sequencen{E_n}$,
adalah saling lepas. \Qed}

\header{212Cd}{\bf (d)}\cmmnt{ Ruang ukur}
$(X,\hat\Sigma,\hat\mu)$ disebut {\bf pelengkapan} dari ruang ukur
$(X,\Sigma,\mu)$; \cmmnt{demikian pula,} saya akan menyebut $\hat\mu$
sebagai {\bf pelengkapan} dari $\mu$, dan kadang-kadang\cmmnt{ (jika
tampak jelas ideal nol mana yang sedang dibicarakan)} saya akan menyebut
$\hat\Sigma$ sebagai pelengkapan dari $\Sigma$. Anggota-anggota
$\hat\Sigma$ kadang-kadang disebut {\bf terukur terhadap $\mu$}.

\leader{212D}{}\cmmnt{Ada sesuatu yang sebaiknya segera saya periksa.

\medskip

\noindent}{\bf Proposisi} Misalkan $(X,\Sigma,\mu)$ sembarang ruang
ukur. Maka $(X,\hat\Sigma,\hat\mu)$, sebagaimana didefinisikan dalam
212C, adalah ruang ukur lengkap dan $\hat\mu$ merupakan perluasan dari
$\mu$; serta $(X,\hat\Sigma,\hat\mu)=(X,\Sigma,\mu)$ jika dan hanya jika
$(X,\Sigma,\mu)$ lengkap.

\proof{{\bf (a)} Misalkan $A\subseteq E\in\hat\Sigma$ dan
$\hat\mu E=0$. Maka (menurut 212Cb) terdapat $E''\in\Sigma$ sedemikian
sehingga $E\subseteq E''$ dan $\mu E''=0$. Dengan demikian kita
mempunyai

\Centerline{$\emptyset\subseteq A\subseteq E''$,\quad
$\mu(E''\setminus\emptyset)=0$,}

\noindent sehingga $A\in\hat\Sigma$. Karena $A$ sembarang, $\hat\mu$
lengkap.

\medskip

{\bf (b)} Jika $E\in\Sigma$, maka tentu saja $E\in\hat\Sigma$, karena
$E\subseteq E\subseteq E$ dan $\mu(E\setminus E)=0$; serta $\hat\mu
E=\mu^*E=\mu E$. Jadi $\Sigma\subseteq\hat\Sigma$ dan $\hat\mu$
memperluas $\mu$.

\medskip

{\bf (c)} Jika $\mu=\hat\mu$, maka tentu saja $\mu$ harus lengkap.
Jika $\mu$ lengkap dan $E\in\hat\Sigma$, maka terdapat $E'$,
$E''\in\Sigma$ sedemikian sehingga $E'\subseteq E\subseteq E''$ dan
$\mu(E''\setminus E')=0$. Tetapi sekarang
$E\setminus E'\subseteq E''\setminus E'$, sehingga (karena
$(X,\Sigma,\mu)$ lengkap) $E\setminus E'\in\Sigma$ dan
$E=E'\cup(E\setminus E')\in\Sigma$. Karena $E$ sembarang,
$\hat\Sigma\subseteq\Sigma$, maka $\hat\Sigma=\Sigma$ dan
$\mu=\hat\mu$.
}%end of proof of 212D

\vleader{72pt}{212E}{}\cmmnt{Pentingnya konstruksi ini sedemikian rupa
sehingga ada gunanya untuk menguraikan beberapa sifat elementer lebih
lanjut.

\medskip

\noindent}{\bf Proposisi} Misalkan $(X,\Sigma,\mu)$ suatu ruang ukur,
dan $(X,\hat\Sigma,\hat\mu)$ pelengkapannya.

(a) Ukuran luar $\hat\mu^*$ dan $\mu^*$ yang didefinisikan dari
$\hat\mu$ dan $\mu$ adalah sama.

(b) $\mu$ dan $\hat\mu$ menghasilkan himpunan-himpunan terabaikan dan
koterabaikan yang sama, serta himpunan-himpunan berukuran luar penuh
yang sama.

(c) $\hat\mu$ adalah satu-satunya ukuran berdomain $\hat\Sigma$ yang
sama dengan $\mu$ pada $\Sigma$.

(d) Suatu subhimpunan dari $X$ termasuk dalam $\hat\Sigma$ jika dan hanya
jika dapat dinyatakan sebagai $F\symmdiff A$, dengan $F\in\Sigma$ dan
$A$ terabaikan terhadap $\mu$.

\proof{{\bf (a)} Ambil sembarang $A\subseteq X$. (i) Jika
$A\subseteq F\in\Sigma$, maka $F\in\hat\Sigma$ dan
$\mu F=\hat\mu F$, sehingga

\Centerline{$\hat\mu^*A\le\hat\mu F=\mu F$;}

\noindent karena $F$ sembarang, $\hat\mu^*A\le\mu^*A$. (ii) Jika
$A\subseteq E\in\hat\Sigma$, terdapat $E''\in\Sigma$ sedemikian
sehingga $E\subseteq E''$ dan $\mu E''=\hat\mu E$, sehingga

\Centerline{$\mu^*A\le\mu E''=\hat\mu E$;}

\noindent karena $E$ sembarang, $\mu^*A\le\hat\mu^*A$.

\medskip
{\bf (b)} Sekarang, untuk $A\subseteq X$,

\Centerline{$A$ terabaikan terhadap $\mu$ jika dan hanya jika
$\mu^*A=0\iff\hat\mu^*A=0\iff A$ terabaikan terhadap $\hat\mu$.}

$$\eqalign{A\text{ koterabaikan terhadap }\mu
&\iff\mu^*(X\setminus A)=0\cr
&\iff\,\hat\mu^*(X\setminus A)=0
\iff A\text{ koterabaikan terhadap }\hat\mu.\cr}$$

\noindent Jika $A$ berukuran luar penuh untuk $\mu$, $F\in\hat\Sigma$,
dan $F\cap A=\emptyset$, maka terdapat $F'\in\Sigma$ sedemikian
sehingga $F'\subseteq F$ dan $\mu F'=\hat\mu F$; karena
$F'\cap A=\emptyset$, $F'$ terabaikan terhadap $\mu$ dan $F$
terabaikan terhadap $\hat\mu$; karena $F$ sembarang, $A$ berukuran luar penuh
untuk $\hat\mu$. Dalam arah sebaliknya, tentu saja, jika $A$ berukuran
luar penuh untuk $\hat\mu$, maka

\Centerline{$\mu^*(F\cap A)=\hat\mu^*(F\cap A)=\hat\mu F=\mu F$}

\noindent untuk setiap $F\in\Sigma$, sehingga $A$ berukuran luar penuh
untuk $\mu$.

\medskip

{\bf (c)} Jika $\tilde\mu$ adalah sembarang ukuran berdomain
$\hat\Sigma$ yang memperluas $\mu$, kita harus mempunyai

\Centerline{$\mu E'\le\tilde\mu E\le\mu E''$,
\quad $\mu E'=\hat\mu E=\mu E''$,}

\noindent sehingga $\tilde\mu E=\hat\mu E$, bilamana
$E'$, $E''\in\Sigma$, $E'\subseteq E\subseteq E''$, dan
$\mu(E''\setminus E')=0$.

\medskip

{\bf (d)(i)} Jika $E\in\hat\Sigma$, ambil $E'$, $E''\in\Sigma$
sedemikian sehingga $E'\subseteq E\subseteq E''$ dan
$\mu(E''\setminus E')=0$. Maka
$E\setminus E'\subseteq E''\setminus E'$, sehingga $E\setminus E'$
terabaikan terhadap $\mu$, dan $E=E'\symmdiff(E\setminus E')$ merupakan selisih
simetris antara suatu anggota $\Sigma$ dan suatu himpunan terabaikan.

\medskip

\quad{\bf (ii)} Jika $E=F\symmdiff A$, dengan $F\in\Sigma$ dan $A$
terabaikan terhadap $\mu$, ambil $G\in\Sigma$ sedemikian sehingga $\mu G=0$
dan $A\subseteq G$; maka $F\setminus G\subseteq E\subseteq F\cup G$
dan $\mu((F\cup G)\setminus(F\setminus G))=\mu G=0$, sehingga
$E\in\hat\Sigma$.
}%end of proof of 212E

\vleader{60pt}{212F}{}\cmmnt{Sekarang mari kita tinjau pengintegralan
terhadap pelengkapan suatu ukuran.

\medskip

\noindent}{\bf Proposisi} Misalkan $(X,\Sigma,\mu)$ suatu ruang ukur
dan $(X,\hat\Sigma,\hat\mu)$ pelengkapannya.

(a) Suatu fungsi bernilai $[-\infty,\infty]$, $f$, yang didefinisikan
pada suatu subhimpunan dari $X$ bersifat terukur terhadap $\hat\Sigma$ jika dan hanya
jika terukur secara virtual terhadap $\mu$.

(b) Misalkan $f$ suatu fungsi bernilai $[-\infty,\infty]$ yang
didefinisikan pada suatu subhimpunan dari $X$. Maka
$\int fd\mu=\int fd\hat\mu$ jika salah satunya didefinisikan dalam
$[-\infty,\infty]$; khususnya, $f$ dapat diintegralkan terhadap $\mu$ jika dan hanya
jika dapat diintegralkan terhadap $\hat\mu$.

\proof{{\bf (a)(i)} Misalkan $f$ suatu fungsi bernilai
$[-\infty,\infty]$ yang terukur terhadap $\hat\Sigma$. Untuk $q\in\Bbb Q$,
misalkan $E_q\in\hat\Sigma$ sedemikian sehingga
$\{x:f(x)\le q\}=\dom f\cap E_q$, dan pilih
$E_q'$, $E_q''\in\Sigma$ sedemikian sehingga
$E_q'\subseteq E_q\subseteq E_q''$ dan
$\mu(E_q''\setminus E_q')=0$. Tetapkan
$H=X\setminus\bigcup_{q\in\Bbb Q}(E_q''\setminus E_q')$; maka $H$
koterabaikan terhadap $\mu$. Untuk $a\in\Bbb R$, tetapkan

\Centerline{$G_a=\bigcup_{q\in\Bbb Q,q<a}E_q'\in\Sigma$;}

\noindent maka

\Centerline{$\{x:x\in\dom(f\restr H),\,(f\restr H)(x)<a\}
=G_a\cap\dom(f\restr H)$.}

\noindent Hal ini menunjukkan bahwa $f\restr H$ terukur terhadap $\Sigma$,
sehingga $f$ terukur secara virtual terhadap $\mu$.

\medskip

\quad{\bf (ii)} Jika $f$ terukur secara virtual terhadap $\mu$, maka
terdapat himpunan $H\subseteq X$ yang koterabaikan terhadap $\mu$ sedemikian sehingga
$f\restr H$ terukur terhadap $\Sigma$. Karena $\Sigma\subseteq\hat\Sigma$,
$f\restr H$ juga terukur terhadap $\hat\Sigma$. Dan $H$
koterabaikan terhadap $\hat\mu$, menurut 212Eb. Namun ini berarti bahwa $f$
terukur secara virtual terhadap $\hat\mu$, sehingga
terukur terhadap $\hat\Sigma$, menurut 212Bb.

\medskip

{\bf (b)(i)} Misalkan $f:D\to[-\infty,\infty]$ suatu fungsi, dengan
$D\subseteq X$. Jika salah satu dari $\int fd\mu$ dan
$\int fd\hat\mu$ didefinisikan dalam $[-\infty,\infty]$, maka $f$
terukur secara virtual dan didefinisikan hampir di mana-mana untuk salah
satu ukuran yang bersangkutan, dan karenanya untuk keduanya (dengan
menggabungkan (a) di atas dengan 212Bb).

\medskip

\quad{\bf (ii)} Sekarang misalkan $f$ nonnegatif dan terintegralkan
terhadap $\mu$ atau terhadap $\hat\mu$. Misalkan $E\in\Sigma$
suatu himpunan koterabaikan yang termuat dalam $\dom f$ sedemikian
sehingga $f\restr E$ terukur terhadap $\Sigma$. Untuk $n\in\Bbb N$, $k\ge 1$,
tetapkan

\Centerline{$E_{nk}=\{x:x\in E$, $f(x)\ge 2^{-n}k\}$;}

\noindent maka setiap $E_{nk}$ termasuk dalam $\Sigma$ dan berukuran
hingga untuk $\mu$ maupun $\hat\mu$. (Jika $f$ dapat diintegralkan terhadap $\mu$,

\Centerline{$\hat\mu E_{nk}=\mu E_{nk}\le 2^n\int fd\mu$;}

\noindent jika $f$ dapat diintegralkan terhadap $\hat\mu$,

\Centerline{$\mu E_{nk}=\hat\mu E_{nk}\le 2^n\int fd\hat\mu$.)}

\noindent Jadi

\Centerline{$f_n=\sum_{k=1}^{4^n}2^{-n}\chi E_{nk}$}

\noindent adalah sederhana terhadap $\mu$ sekaligus sederhana terhadap $\hat\mu$, dan
$\int f_nd\mu=\int f_nd\hat\mu$. Perhatikan bahwa, untuk $x\in E$,

$$\eqalign{f_n(x)
&=2^{-n}k\text{ jika }k<4^n\text{ dan }2^{-n}k\le f(x)<2^{-n}(k+1),\cr
&=2^n\text{ jika }f(x)\ge 2^n.\cr}$$

\noindent Jadi $\sequencen{f_n}$ merupakan barisan fungsi tak menurun
yang konvergen ke $f$ di setiap titik $E$, yakni, hampir di mana-mana
terhadap $\mu$ maupun terhadap $\hat\mu$. Karena itu, untuk setiap
$c\in\Bbb R$, kita mempunyai

$$\eqalign{\int fd\mu=c
&\iff\lim_{n\to\infty}\int f_nd\mu=c\cr
&\iff\lim_{n\to\infty}\int f_nd\hat\mu=c
\iff\int fd\hat\mu=c.\cr}$$

\medskip

\quad{\bf (iii)} Mengenai integral tak hingga, ingat bahwa untuk fungsi
nonnegatif saya menuliskan `$\int f=\infty$' tepat ketika $f$
didefinisikan hampir di mana-mana, terukur secara virtual, dan tidak
terintegralkan. Jadi (i) dan (ii) bersama-sama menunjukkan bahwa
$\int fd\mu=\int fd\hat\mu$ bilamana $f$ nonnegatif dan salah satu
integral didefinisikan dalam $[0,\infty]$.

\medskip

\quad{\bf (iv)} Karena $\mu$ maupun $\hat\mu$ sama-sama menetapkan bahwa
$\int f$ harus ditafsirkan sebagai $\int f^+-\int f^-$ tepat ketika
ungkapan ini dapat didefinisikan dalam $[-\infty,\infty]$, dengan
menuliskan $f^+(x)=\max(f(x),0)$ dan $f^-(x)=\max(-f(x),0)$ untuk
$x\in\dom f$, hasil untuk $f$ bernilai real umum segera diperoleh.
}%end of proof of 212F

\leader{212G}{}\cmmnt{Sekarang saya beralih ke persoalan tentang
pengaruh konstruksi ini terhadap sifat-sifat yang tercantum dalam
211B-211K. %211B 211C 211D 211E 211F 211G 211H 211I 211J 211K
\medskip

\noindent}{\bf Proposisi} Misalkan $(X,\Sigma,\mu)$ suatu ruang ukur,
dan $(X,\hat\Sigma,\hat\mu)$ pelengkapannya.

(a) Untuk masing-masing sifat berikut -- ruang peluang, berhingga total,
$\sigma$-hingga, semihingga, dan dapat dilokalkan --
$(X,\hat\Sigma,\hat\mu)$ memiliki sifat tersebut jika dan hanya jika
$(X,\Sigma,\mu)$ memilikinya.

(b) $(X,\hat\Sigma,\hat\mu)$ dapat dilokalkan secara ketat jika
$(X,\Sigma,\mu)$ demikian, dan setiap dekomposisi $X$ untuk $\mu$
merupakan dekomposisi untuk $\hat\mu$.

(c) Suatu himpunan $H\in\hat\Sigma$ merupakan atom untuk $\hat\mu$ jika
dan hanya jika terdapat $E\in\Sigma$ sedemikian sehingga $E$ merupakan
atom untuk $\mu$ dan $\hat\mu(H\symmdiff E)=0$.

(d) $(X,\hat\Sigma,\hat\mu)$ tanpa atom atau atomik murni jika dan hanya
jika $(X,\Sigma,\mu)$ demikian.

\proof{{\bf (a)(i)} Karena $\hat\mu X=\mu X$,
$(X,\hat\Sigma,\hat\mu)$ merupakan ruang peluang, atau hingga total,
jika dan hanya jika $(X,\Sigma,\mu)$ demikian.

\medskip

\quad{\bf (ii)}\grheada\ Jika $(X,\Sigma,\mu)$ bersifat $\sigma$-hingga,
terdapat suatu barisan $\sequencen{E_n}$ yang menutupi $X$, dengan
$\mu E_n<\infty$ untuk setiap $n$. Sekarang $\hat\mu E_n<\infty$ untuk
setiap $n$, sehingga $(X,\hat\Sigma,\hat\mu)$ bersifat $\sigma$-hingga.

\medskip

\qquad\grheadb\ Jika $(X,\hat\Sigma,\hat\mu)$ bersifat $\sigma$-hingga,
terdapat suatu barisan $\sequencen{E_n}$ yang menutupi $X$, dengan
$\hat\mu E_n<\infty$ untuk setiap $n$. Sekarang untuk setiap $n$ kita
dapat menemukan $E_n''\in\Sigma$ sedemikian sehingga
$\mu E_n''<\infty$ dan $E_n\subseteq E_n''$; jadi
$\sequencen{E_n''}$ menyaksikan bahwa $(X,\Sigma,\mu)$ bersifat
$\sigma$-hingga.

\medskip

\quad{\bf (iii)}\grheada\ Jika $(X,\Sigma,\mu)$ semihingga dan
$\hat\mu E=\infty$, maka terdapat $E'\in\Sigma$ sedemikian sehingga
$E'\subseteq E$ dan $\mu E'=\infty$. Selanjutnya, terdapat
$F\in\Sigma$ sedemikian sehingga $F\subseteq E'$ dan
$0<\mu F<\infty$. Tentu saja sekarang kita mempunyai
$F\in\hat\Sigma$, $F\subseteq E$, dan $0<\hat\mu F<\infty$. Karena
$E$ sembarang, $(X,\hat\Sigma,\hat\mu)$ semihingga.

\medskip

\qquad\grheadb\ Jika $(X,\hat\Sigma,\hat\mu)$ semihingga dan
$\mu E=\infty$, maka $\hat\mu E=\infty$, sehingga terdapat
$F\subseteq E$ sedemikian sehingga $0<\hat\mu F<\infty$. Selanjutnya,
terdapat $F'\in\Sigma$ sedemikian sehingga $F'\subseteq F$ dan
$\mu F'=\hat\mu F$. Tentu saja sekarang kita mempunyai
$F'\subseteq E$ dan $0<\mu F'<\infty$. Karena $E$ sembarang,
$(X,\Sigma,\mu)$ semihingga.

\medskip

\quad{\bf (iv)}\grheada\ Jika $(X,\Sigma,\mu)$ dapat dilokalkan dan
$\Cal E\subseteq \hat\Sigma$, maka tetapkan

\Centerline{$\Cal F=\{F:F\in\Sigma,\,\exists\enskip
E\in\Cal E,\, F\subseteq E\}$.}

\noindent Misalkan $H$ suatu supremum esensial dari $\Cal F$ dalam
$\Sigma$, seperti dalam 211G.

Jika $E\in\Cal E$, terdapat $E'\in\Sigma$ sedemikian sehingga
$E'\subseteq E$ dan $E\setminus E'$ terabaikan; sekarang
$E'\in\Cal F$, sehingga

\Centerline{$\hat\mu(E\setminus H)
\le\hat\mu(E\setminus E')+\mu(E'\setminus H)=0$.}

\noindent Jika $G\in\hat\Sigma$ dan $\hat\mu(E\setminus G)=0$ untuk
setiap $E\in\Cal E$, misalkan $G''\in\Sigma$ sedemikian sehingga
$G\subseteq G''$ dan $\hat\mu(G''\setminus G)=0$; kemudian, untuk setiap
$F\in\Cal F$, terdapat $E\in\Cal E$ yang memuat $F$, sehingga

\Centerline{$\mu(F\setminus G'')\le\hat\mu(E\setminus G)=0.$}

\noindent Karena $F$ sembarang, $\mu(H\setminus G'')=0$ dan
$\hat\mu(H\setminus G)=0$. Hal ini menunjukkan bahwa $H$ merupakan
supremum esensial dari $\Cal E$ dalam $\hat\Sigma$. Karena $\Cal E$
sembarang, $(X,\hat\Sigma,\hat\mu)$ dapat dilokalkan.

\medskip

\qquad\grheadb\ Misalkan $(X,\hat\Sigma,\hat\mu)$ dapat dilokalkan dan
$\Cal E\subseteq\Sigma$. Dengan bekerja dalam
$(X,\hat\Sigma,\hat\mu)$, misalkan $H$ suatu supremum esensial untuk
$\Cal E$ dalam $\hat\Sigma$. Misalkan $H'\in\Sigma$ sedemikian sehingga
$H'\subseteq H$ dan $\hat\mu(H\setminus H')=0$. Maka

\Centerline{$\mu(E\setminus H')
\le\hat\mu(E\setminus H)+\hat\mu(H\setminus H')=0$}

\noindent untuk setiap $E\in\Cal E$; sedangkan jika $G\in\Sigma$ dan
$\mu(E\setminus G)=0$ untuk setiap $E\in\Cal E$, kita harus mempunyai

\Centerline{$\mu(H'\setminus G)\le\hat\mu(H\setminus G)=0$.}

\noindent Jadi $H'$ merupakan supremum esensial dari $\Cal E$ dalam
$\Sigma$. Karena $\Cal E$ sembarang, $(X,\Sigma,\mu)$ dapat dilokalkan.

\medskip

{\bf (b)} Misalkan $\langle X_i\rangle_{i\in I}$ suatu dekomposisi $X$
untuk $\mu$, seperti dalam 211E. Tentu saja ini merupakan partisi $X$
menjadi himpunan-himpunan berukuran hingga menurut $\hat\mu$. Jika $H\subseteq X$
dan $H\cap X_i\in\hat\Sigma$ untuk setiap $i$, untuk setiap $i\in I$
pilih himpunan $E_i'$, $E_i''\in\Sigma$ sedemikian sehingga

\Centerline{$E_i'\subseteq H\cap X_i\subseteq E_i''$,\quad
$\mu(E_i''\setminus E_i')=0$.}

\noindent Tetapkan $E'=\bigcup_{i\in I}E_i'$,
$E''=\bigcup_{i\in I}(E_i''\cap X_i)$. Maka $E'\cap X_i=E_i'$,
$E''\cap X_i=E_i''\cap X_i$ untuk setiap $i$, sehingga $E'$ dan $E''$
termasuk dalam $\Sigma$. Juga

\Centerline{$\mu(E''\setminus E')
=\sum_{i\in I}\mu(E_i''\cap X_i\setminus E_i')=0$.}

\noindent Karena $E'\subseteq H\subseteq E''$, $H\in\hat\Sigma$ dan

\Centerline{$\hat\mu H=\mu E'
=\sum_{i\in I}\mu E'_i=\sum_{i\in I}\hat\mu(H\cap X_i)$.}

\noindent Karena $H$ sembarang, $\langle X_i\rangle_{i\in I}$
merupakan dekomposisi $X$ untuk $\hat\mu$.

Dengan demikian, $(X,\hat\Sigma,\hat\mu)$ dapat dilokalkan secara ketat
jika dekomposisi semacam itu ada, dan hal ini berlaku jika
$(X,\Sigma,\mu)$ dapat dilokalkan secara ketat.

\medskip

{\bf (c)-(d)(i)} Misalkan $E\in\hat\Sigma$ suatu atom untuk $\hat\mu$.
Misalkan $E'\in\Sigma$ sedemikian sehingga $E'\subseteq E$ dan
$\hat\mu(E\setminus E')=0$. Maka $\mu E'=\hat\mu E>0$. Jika
$F\in\Sigma$ dan $F\subseteq E'$, maka $F\subseteq E$, sehingga
$\mu F=\hat\mu F=0$ atau $\mu(E'\setminus F)=\hat\mu(E\setminus F)=0$.
Karena $F$ sembarang, $E'$ merupakan atom untuk $\mu$, dan
$\hat\mu(E\symmdiff E')=\hat\mu(E\setminus E')=0$.

\medskip

\quad{\bf (ii)} Misalkan $E\in\Sigma$ suatu atom untuk $\mu$, dan
$H\in\hat\Sigma$ sedemikian sehingga $\hat\mu(H\symmdiff E)=0$. Maka
$\hat\mu H=\mu E>0$. Jika $F\in\hat\Sigma$ dan $F\subseteq H$,
misalkan $F'\subseteq F$ sedemikian sehingga $F'\in\Sigma$ dan
$\hat\mu(F\setminus F')=0$. Maka $E\cap F'\subseteq E$ dan
$\hat\mu(F\symmdiff(E\cap F'))=0$, sehingga $\hat\mu F=\mu(E\cap F')=0$
atau $\hat\mu(H\setminus F)=\mu(E\setminus F')=0$. Karena $F$
sembarang, $H$ merupakan atom untuk $\hat\mu$.

\medskip

\quad{\bf (iii)} Segera diperoleh bahwa $(X,\hat\Sigma,\hat\mu)$ tanpa
atom jika dan hanya jika $(X,\Sigma,\mu)$ tanpa atom.

\medskip

\quad{\bf (iv)}\grheada\ Di pihak lain, jika $(X,\Sigma,\mu)$ atomik
murni dan $\hat\mu H>0$, terdapat $E\in\Sigma$ sedemikian sehingga
$E\subseteq H$ dan $\mu E>0$, serta suatu atom $F$ untuk $\mu$
sedemikian sehingga $F\subseteq E$; tetapi $F$ juga merupakan atom untuk
$\hat\mu$. Karena $H$ sembarang, $(X,\hat\Sigma,\hat\mu)$ atomik murni.

\medskip

\qquad\grheadb\ Dan jika $(X,\hat\Sigma,\hat\mu)$ atomik murni dan
$\mu E>0$, maka terdapat $H\subseteq E$ yang merupakan atom untuk
$\hat\mu$; sekarang misalkan $F\in\Sigma$ sedemikian sehingga
$F\subseteq H$ dan $\hat\mu(H\setminus F)=0$, sehingga $F$ merupakan
atom untuk $\mu$ dan $F\subseteq E$. Karena $E$ sembarang,
$(X,\Sigma,\mu)$ atomik murni.
}%end of proof of 212G

\exercises{
\leader{212X}{Latihan dasar $\pmb{>}$(a)} 
%\spheader 212Xa
Misalkan $(X,\Sigma,\mu)$ suatu ruang ukur lengkap. Andaikan
$A\subseteq E\in\Sigma$ dan
$\mu^*A+\mu^*(E\setminus A)=\mu E<\infty$. Tunjukkan bahwa
$A\in\Sigma$.

\sqheader 212Xb Misalkan $\mu$ dan $\nu$ dua ukuran pada suatu himpunan
$X$, dengan pelengkapan $\hat\mu$ dan $\hat\nu$. Tunjukkan bahwa
pernyataan-pernyataan berikut ekuivalen: (i) ukuran luar $\mu^*$ dan
$\nu^*$ yang didefinisikan dari $\mu$ dan $\nu$ adalah sama;
(ii) $\hat\mu E=\hat\nu E$ bilamana salah satunya didefinisikan dan
hingga; (iii) $\int fd\mu=\int fd\nu$ bilamana $f$ suatu fungsi bernilai
real sedemikian sehingga salah satu integral didefinisikan dan hingga.
\Hint{untuk (i)$\Rightarrow$(ii), jika $\hat\mu E<\infty$, ambil suatu
selubung terukur $F$ dari $E$ untuk $\nu$ dan hitung
$\nu^*E+\nu^*(F\setminus E)$.}

\spheader 212Xc Misalkan $\mu$ adalah pembatasan ukuran Lebesgue pada
aljabar-sigma Borel dari $\Bbb R$, seperti dalam 211P. Tunjukkan bahwa
pelengkapannya adalah ukuran Lebesgue itu sendiri. \Hint{134F.}

\spheader 212Xd Ulangi 212Xc untuk (i) ukuran Lebesgue pada
$\BbbR^r$ (ii) ukuran Lebesgue--Stieltjes pada $\Bbb R$ (114Xa).

\spheader 212Xe
Misalkan $X$ suatu himpunan dan $\Sigma$ suatu aljabar-sigma dari
subhimpunan-subhimpunan $X$. Misalkan $\Cal I$ suatu ideal-$\sigma$ dari
subhimpunan-subhimpunan $X$ (112Db).
(i) Tunjukkan bahwa $\Sigma_1=\{E\symmdiff A:E\in\Sigma,\,A\in\Cal I\}$
adalah suatu aljabar-sigma dari subhimpunan-subhimpunan $X$.
(ii) Misalkan $\Sigma_2$ keluarga himpunan $E\subseteq X$ sedemikian
sehingga terdapat $E'$, $E''\in\Sigma$ dengan
$E'\subseteq E\subseteq E''$ dan $E''\setminus E'\in\Cal I$.
Tunjukkan bahwa $\Sigma_2$ merupakan suatu aljabar-sigma dari
subhimpunan-subhimpunan $X$ dan bahwa $\Sigma_2\subseteq\Sigma_1$.
(iii) Tunjukkan bahwa $\Sigma_2=\Sigma_1$ jika dan hanya jika setiap
anggota $\Cal I$ termuat dalam suatu anggota $\Sigma\cap\Cal I$.

\spheader 212Xf
Misalkan $(X,\Sigma,\mu)$ suatu ruang ukur, $Y$ sembarang himpunan, dan
$\phi:X\to Y$ suatu fungsi. Tetapkan $\theta B=\mu^*\phi^{-1}[B]$ untuk
setiap $B\subseteq Y$. (i) Tunjukkan bahwa $\theta$ adalah ukuran luar
pada $Y$. (ii) Misalkan $\nu$ ukuran yang didefinisikan dari $\theta$
dengan metode \Caratheodory, dan $\Tau$ domainnya. Tunjukkan bahwa jika
$C\subseteq Y$ dan $\phi^{-1}[C]\in\Sigma$, maka $C\in\Tau$.
(iii) Andaikan $(X,\Sigma,\mu)$ lengkap dan hingga total. Tunjukkan
bahwa $\nu$ adalah ukuran citra $\mu\phi^{-1}$.

\spheader 212Xg
Misalkan $g$ dan $h$ dua fungsi tak menurun dari $\Bbb R$ ke dirinya
sendiri, dan $\mu_g$, $\mu_h$ ukuran Lebesgue--Stieltjes yang terkait.
Tunjukkan bahwa suatu fungsi bernilai real $f$ yang didefinisikan pada
suatu subhimpunan dari $\Bbb R$ dapat diintegralkan terhadap $\mu_{g+h}$ jika dan hanya
jika fungsi tersebut dapat diintegralkan terhadap $\mu_g$ sekaligus
dapat diintegralkan terhadap $\mu_h$, dan bahwa dalam hal itu
$\int fd\mu_{g+h}=\int fd\mu_g+\int fd\mu_h$. \Hint{114Yb}.

\spheader 212Xh Misalkan $(X,\Sigma,\mu)$ suatu ruang ukur, dan
$\Cal I$ suatu ideal-$\sigma$ dari subhimpunan-subhimpunan $X$; tetapkan
$\Sigma_1=\{E\symmdiff A:E\in\Sigma,\,A\in\Cal I\}$, seperti dalam
212Xe. Tunjukkan bahwa jika setiap anggota $\Sigma\cap\Cal I$
terabaikan terhadap $\mu$, maka terdapat perluasan unik dari $\mu$ menjadi suatu
ukuran $\mu_1$ berdomain $\Sigma_1$ sedemikian sehingga $\mu_1A=0$
untuk setiap $A\in\Cal I$.
%212C

\spheader 212Xi\dvAnew{2009}
Misalkan $(X,\Sigma,\mu)$ suatu ruang ukur lengkap sedemikian sehingga
$\mu X>0$, $Y$ suatu himpunan, $f:X\to Y$ suatu fungsi, dan
$\mu f^{-1}$ ukuran citra pada $Y$. Tunjukkan bahwa jika $\Cal F$
adalah filter subhimpunan-subhimpunan koterabaikan terhadap $\mu$ dari $X$, maka
filter citra $f[[\Cal F]]$ (2A1Ib) adalah filter subhimpunan-subhimpunan
koterabaikan terhadap $\mu f^{-1}$ dari $Y$.

\spheader 212Xj\dvAnew{2010}
Misalkan $(X,\Sigma,\mu)$ suatu ruang ukur lengkap dan $f:X\to\Bbb R$
suatu fungsi sedemikian sehingga $\overline{\int}fd\mu<\infty$.
Tunjukkan bahwa terdapat fungsi terukur $g:X\to\Bbb R$ sedemikian
sehingga $f(x)\le g(x)$ untuk setiap $x\in X$ dan
$\int g\,d\mu=\overline{\int}fd\mu$.

\leader{212Y}{Latihan lanjutan (a)}
%\spheader 212Ya
Misalkan $X$ suatu himpunan dan $\phi$ suatu ukuran dalam pada $X$,
yakni, suatu fungsional dari $\Cal PX$ ke $[0,\infty]$ sedemikian
sehingga

\inset{$\phi\emptyset=0$,}
\inset{$\phi(A\cup B)\ge\phi A+\phi B$ jika $A\cap B=\emptyset$,}

\inset{$\phi(\bigcap_{n\in\Bbb N}A_n)=\lim_{n\to\infty}\phi A_n$
bilamana $\sequencen{A_n}$ adalah suatu barisan tak menaik dari
subhimpunan-subhimpunan $X$ dan $\phi A_0<\infty$,}

\inset{jika $\phi A=\infty$, $a\in\Bbb R$, terdapat $B\subseteq A$
sedemikian sehingga $a\le\phi B<\infty$.}

\noindent Misalkan $\mu$ ukuran yang didefinisikan dari $\phi$, yakni,
$\mu=\phi\restr\Sigma$, dengan

\Centerline{$\Sigma=\{E:\phi(A)
  =\phi(A\cap E)+\phi(A\setminus E)\Forall A\subseteq X\}$}

\noindent (113Yg). Tunjukkan bahwa $\mu$ harus lengkap.
}%end of exercises

\endnotes{
\Notesheader{212} Proses pelengkapan sedemikian alami, dapat diterapkan
secara universal, dan praktis, sehingga dalam bagian-bagian besar teori
ukuran adalah wajar untuk hanya menggunakan ruang ukur lengkap. Bahkan,
banyak penulis merumuskan definisi mereka sedemikian rupa sehingga,
secara eksplisit ataupun implisit, hanya ruang ukur lengkap yang
dipertimbangkan. Dalam risalah ini saya menghindari mengambil langkah
sebesar itu, meskipun langkah tersebut akan menyederhanakan pernyataan
banyak teorema dalam jilid ini (misalnya). Dalam Jilid 1 saya memang
bersusah payah memberikan definisi `fungsi terintegralkan' yang, pada
hakikatnya, menelaah keterintegralan terhadap pelengkapan suatu ukuran
(212Fb). Terdapat ruang ukur tidak lengkap yang layak dikaji (misalnya,
pembatasan ukuran Lebesgue pada aljabar-sigma Borel dari $\Bbb R$ --
lihat 211P), dan beberapa persoalan menarik yang akan dibahas dalam
Jilid 3 dan 5 berlaku bagi ruang-ruang tersebut. Oleh karena itu, dengan
menanggung cukup banyak manuver verbal, saya lebih suka menuliskan
teorema dalam bentuk yang dapat diterapkan pada ruang ukur sembarang,
tanpa mengasumsikan kelengkapan. Namun, akan wajar, dan memang akan
mempertajam teknik Anda, jika Anda secara teratur mencari perumusan
alternatif yang menjadi alami apabila Anda hanya berminat pada ruang
lengkap.
}

\discrpage



